\clearpage
\VerdierPhysicalPageAnchor{0024}
谱序列的机制往往掩盖了现象本身的简洁性。Künneth 关系的初等例子对此颇有
代表性。

\VerdierSemanticAnchor{intro-5}
\noindent\textbf{5.}\quad
设 $X$ 与 $Y$ 是两个紧空间，$A$ 是有单位元的交换环（为简单起见），
$F$ 与 $G$ 分别是 $X$ 与 $Y$ 上的 $A$-模层。当 $A$ 是域时，已知外积
[23] 给出分次向量空间的同构：

\VerdierSemanticAnchor{intro-5-kunneth-field}
\begin{equation}
\mathrm H^*(X,F)\otimes_A\mathrm H^*(Y,G)
\xrightarrow{\sim}
\mathrm H^*(X\times Y,F\otimes_A G).
\tag{$\text{Kü}_{\mathrm{coh}}^{1}$}
\end{equation}

当 $A=\mathbb Z$，并且层 $F$ 或层 $G$ 无挠时，有可分裂正合列：

\VerdierSemanticAnchor{intro-5-kunneth-torsion-free}
\begin{equation}
\begin{aligned}
0 &\longrightarrow
\bigoplus_{p+q=n}\mathrm H^p(X,F)\otimes_A\mathrm H^q(Y,G)
\longrightarrow \mathrm H^n(X\times Y,F\otimes_A G) \longrightarrow \\
&\longrightarrow
\bigoplus_{p+q=n+1}
\operatorname{Tor}^{\mathbb Z}\bigl(\mathrm H^p(X,F),\mathrm H^q(Y,G)\bigr)
\longrightarrow 0.
\end{aligned}
\tag{$\text{Kü}_{\mathrm{coh}}^{2}$}
\end{equation}

当 $A=\mathbb Z$，而不对 $F$ 与 $G$ 作任何假设时，有两族可分裂正合列：

\VerdierSemanticAnchor{intro-5-kunneth-general-split}
\begin{equation}
\begin{aligned}
0 &\longrightarrow
\bigoplus_{p+q=n}\mathrm H^p(X,F)\otimes_A\mathrm H^q(Y,G)
\longrightarrow \mathbb H^n \longrightarrow \\
&\longrightarrow
\bigoplus_{p+q=n+1}
\operatorname{Tor}^{\mathbb Z}\bigl(\mathrm H^p(X,F),\mathrm H^q(Y,G)\bigr)
\longrightarrow 0, \\
0 &\longrightarrow
\mathrm H^{n+1}\bigl(X\times Y,\mathcal T\!or^A(F,G)\bigr)
\longrightarrow \mathbb H^n
\longrightarrow \mathrm H^n(X\times Y,F\otimes_A G)
\longrightarrow 0.
\end{aligned}
\tag{$\text{Kü}_{\mathrm{coh}}^{3}$}
\end{equation}

因此，最初的表述 $\text{(Kü)}_{\mathrm{coh}}^{1}$ 变得复杂，但仍可作
显式计算。当环 $A$ 更复杂时，例如
$A=\mathbb Z/\ell^r\mathbb Z$，其中 $\ell$ 为素数、$r>1$；这是一个全局
上同调维数无限的环。传统方法只给出两个具有同一收敛目标的谱序列，其初始
项为：

\VerdierSemanticAnchor{intro-5-kunneth-spectral-sequences}
\begin{equation}
\begin{aligned}
{}'K_2^{p,q}
  &=\bigoplus_{r+s=q}
    \operatorname{Tor}_{-p}^{A}
    \bigl(\mathrm H^r(X,F),\mathrm H^s(Y,G)\bigr), \\
{}''K_2^{p,q}
  &=\mathrm H^p\bigl(X\times Y,\mathcal T\!or_{-q}^{A}(F,G)\bigr).
\end{aligned}
\tag{$\text{Kü}_{\mathrm{coh}}^{4}$}
\end{equation}

然而，这两个谱序列及其收敛目标之间的同构，只是 $A$-模范畴的导出范畴中
下述关系的推论与不完全表达：

\VerdierSemanticAnchor{intro-5-kunneth-derived}
\begin{equation}
\mathrm{RH}(X,F)\overset{\mathrm L}{\otimes}_A\mathrm{RH}(Y,G)
\xrightarrow{\sim}
\mathrm{RH}\bigl(X\times Y,F\overset{\mathrm L}{\otimes}_A G\bigr).
\tag{$\text{Kü}_{\mathrm{Der}}$}
\end{equation}
